% from ega1-3-fr.tex
\oldpage[I]{118}
En particulier, si $U$ est une partie ouverte de $X$, le préschéma induit sur
$U\cap f^{-1}(y)$ par le préschéma $f^{-1}(y)$ est isomorphe à
$f_U^{-1}(y)$ ($f_U$ étant la restriction de $f$ à $U$).
\end{env}

\begin{proposition}[3.6.4) («~transitivité des fibres~»]
\label{I.3.6.4-fr}
Soient $f:X\to Y$, $g:Y'\to Y$ deux morphismes~; posons
$X'=X\times_Y Y'=X_{(Y')}$, et
$f'=f_{(Y')}:X'\to Y'$. Pour tout $y'\in Y'$, si on pose $y=g(y')$, le
préschéma $f'^{-1}(y')$ est isomorphe à
$f^{-1}(y)\otimes_{k(y)}k(y')$.
\end{proposition}

En effet, cela revient à remarquer que les deux préschémas
$(X\otimes_Y k(y))\otimes_{k(y)}k(y')$ et
$(X\times_Y Y')\otimes_{Y'}k(y')$ sont tous deux canoniquement isomorphes à
$X\times_Y\operatorname{Spec}(k(y'))$ par
(\hyperref[I.3.3.9.1-fr]{3.3.9.1}).

En particulier, si $V$ est un voisinage ouvert de $y$ dans $Y$, et si on
désigne par $f_V$ la restriction de $f$ au préschéma induit sur $f^{-1}(V)$,
les préschémas $f^{-1}(y)$ et $f_V^{-1}(y)$ s'identifient canoniquement.

\begin{proposition}[3.6.5]
\label{I.3.6.5-fr}
Soient $f:X\to Y$ un morphisme, $y$ un point de $Y$, $Z$ le préschéma local
$\operatorname{Spec}(\mathcal{O}_y)$, $p=(\psi,\theta)$ la projection
$X\times_Y Z\to X$~; $p$ est un homéomorphisme de l'espace sous-jacent à
$X\times_Y Z$ sur le sous-espace $f^{-1}(Z)$ de $X$ (lorsque l'espace
sous-jacent à $Z$ est identifié à un sous-espace de $Y$, cf.
(\hyperref[I.2.4.2-fr]{2.4.2})), et pour tout $t\in X\times_Y Z$, si on pose
$z=\psi(t)$, $\theta_t^\sharp$ est un isomorphisme de $\mathcal{O}_z$ sur
$\mathcal{O}_t$.
\end{proposition}

Comme $Z$ (identifié à un sous-espace de $Y$) est contenu dans tout ouvert
affine contenant $y$ (\hyperref[I.2.4.2-fr]{2.4.2}), on peut, comme dans
(\hyperref[I.3.6.1-fr]{3.6.1}) se ramener au cas où
$X=\operatorname{Spec}(A)$ et $Y=\operatorname{Spec}(B)$ sont des schémas
affines, $A$ étant une $B$-algèbre. Alors $X\times_Y Z$ est le spectre premier
de $A\otimes_B B_y$ et cet anneau s'identifie canoniquement à $S^{-1}A$, où
$S$ est l'image de $B-\mathfrak{j}_y$ dans $A$ (0, 1.5.2)~; comme $p$
correspond alors à l'homomorphisme canonique $A\to S^{-1}A$, la proposition
résulte de (\hyperref[I.1.6.2-fr]{1.6.2}).

\subsection{Application : réduction d'un préschéma mod.
\texorpdfstring{$\mathfrak{J}$\protect\footnotemark}{J}}
\label{subsection:I.3.7-fr}
\label{I.3.7-fr}
\footnotetext{Ce numéro, qui fait usage de notions et résultats de la suite du
chap. I\textsuperscript{er} et du chap. II, ne sera pas utilisé par la suite,
et n'est destiné qu'aux lecteurs familiers avec la Géométrie algébrique
classique.}

\begin{env}[3.7.1]
\label{I.3.7.1-fr}
Soient $A$ un anneau, $X$ un $A$-préschéma, $\mathfrak{J}$ un idéal de $A$~;
alors $X_0=X\otimes_A(A/\mathfrak{J})$ est un
$(A/\mathfrak{J})$-préschéma, dont on dit parfois qu'il est déduit de $X$ par
\emph{réduction} mod. $\mathfrak{J}$.
\end{env}

\begin{env}[3.7.2]
\label{I.3.7.2-fr}
Cette terminologie est surtout utilisée lorsque $A$ est un \emph{anneau
local} et $\mathfrak{J}$ son idéal maximal, de sorte que $X_0$ est un
préschéma sur le corps résiduel $k=A/\mathfrak{J}$ de $A$.

Lorsqu'en outre $A$ est intègre, de corps des fractions $K$, on peut aussi
considérer le $K$-préschéma $X'=X\otimes_A K$. Par un abus de langage dont
nous ne ferons pas usage, on disait d'ordinaire jusqu'à présent que $X_0$ est
\emph{déduit de $X'$} par réduction mod. $\mathfrak{J}$. Dans les cas où ce
langage était utilisé, $A$ était un anneau local de dimension 1 (le plus
souvent un anneau de valuation discrète) et il était sous-entendu (de façon
plus ou moins explicite) que le $K$-préschéma $X'$ donné était un
sous-préschéma fermé d'un $K$-préschéma $P'$ (en fait, un espace projectif
type $\mathbf{P}_K^r$, cf. II, 4.1.1), lui-même de la forme
$P'=P\otimes_A K$, où $P$ est un $A$-préschéma donné (en fait, le $A$-schéma
$\mathbf{P}_A^r$, avec les notations de II, 4.1.1). Dans notre langage, la
définition de $X_0$ à partir de $X'$ se formule comme suit~:

On considère le schéma affine $Y=\operatorname{Spec}(A)$, formé de deux
points, l'unique point
